\section{Reinforcement Learning Algorithms}

This work evaluates five standard Deep Reinforcement Learning algorithms, covering both 
on-policy and off-policy families. The goal is to compare how different learning paradigms 
behave under the same trading environment and feature representations, while using carefully 
tuned hyperparameter configurations that ensure stable and fair training across methods.

\subsection{On-Policy Methods}

On-policy methods learn directly from data generated by the current policy, meaning that 
the policy is continuously updated using fresh trajectories collected from interaction with 
the environment. This typically leads to more stable learning dynamics, since the data 
distribution closely matches the policy being optimized. However, this comes at the cost 
of lower sample efficiency, as past experiences cannot be reused.

\paragraph{Advantage Actor-Critic (A2C).} 
A2C is a synchronous actor-critic algorithm that jointly learns a policy function (actor) 
and a value function (critic). The critic estimates the expected return, while the actor 
updates the policy using the advantage signal, reducing variance compared to standard policy 
gradient methods.

A2C uses a rollout length of \texttt{n\_steps = 128}, which controls how many
environment steps are collected before each update. Smaller values lead to more frequent but
noisier updates, while larger values improve gradient stability at the cost of responsiveness.

The discount factor \texttt{gamma = 0.99} determines how much the agent prioritizes long-term
rewards, which is essential in financial settings where gains are often delayed. The parameter
\texttt{gae\_lambda = 0.95} controls the bias-variance trade-off in advantage estimation,
with lower values favoring bias reduction and higher values improving variance reduction.

Exploration is encouraged via \texttt{ent\_coef = 0.01}, which prevents premature convergence
to deterministic policies, while \texttt{vf\_coef = 0.5} balances policy learning with value
function accuracy. Gradient stability is further ensured using \texttt{max\_grad\_norm = 0.5}
to prevent excessively large updates.

\paragraph{Proximal Policy Optimization (PPO).} 
PPO improves upon standard policy gradient methods by constraining policy updates through a 
clipped objective, preventing large and unstable policy shifts.

A large rollout size of \texttt{n\_steps = 2048} is used to ensure low-variance advantage
estimation before each optimization phase. The optimization itself is performed using
mini-batches of size \texttt{batch\_size = 64}, which provides a balance between gradient
stability and computational efficiency.

The parameter \texttt{target\_kl = 0.02} plays a key role in stabilizing training by limiting
the Kullback-Leibler divergence between successive policies. Lower values enforce more conservative
updates, while higher values allow faster but potentially unstable learning.

As in A2C, the entropy coefficient \texttt{ent\_coef = 0.01} is used to maintain sufficient
exploration during training, which is particularly important in non-stationary financial
environments.

\subsection{Off-Policy Methods}

Off-policy methods learn from experiences stored in a replay buffer, allowing the agent 
to reuse past transitions generated by previous policies. This significantly improves 
sample efficiency and enables more stable learning in environments where data collection 
is expensive or limited. However, these methods are generally more sensitive to hyperparameter
choices and function approximation errors.

\paragraph{Deep Deterministic Policy Gradient (DDPG).} 
DDPG is an off-policy actor-critic method for continuous action spaces that learns a 
deterministic policy.

A replay buffer of size \texttt{buffer\_size = 20{,}000} stores past transitions, enabling
learning from previously experienced market conditions. The \texttt{batch\_size = 256}
controls the number of samples used per update, where larger values improve gradient stability
but increase computational cost.

Training begins after an initial exploration phase of \texttt{learning\_starts = 1000} steps,
ensuring that the replay buffer contains sufficiently diverse experiences before learning
commences. The learning rate \texttt{1e-5} is intentionally small to ensure stable updates
in the presence of highly noisy financial signals.

\paragraph{Twin Delayed DDPG (TD3).} 
TD3 extends DDPG by introducing twin critic networks, delayed policy updates, and target policy 
smoothing, significantly improving stability.

It shares the same core hyperparameters as DDPG, including \texttt{buffer\_size = 20{,}000},
\texttt{batch\_size = 256}, and \texttt{learning\_starts = 1000}. These settings ensure a
consistent comparison with DDPG while benefiting from TD3’s architectural improvements that
reduce overestimation bias and improve value function reliability.

\paragraph{Soft Actor-Critic (SAC).} 
SAC is an off-policy method that optimizes a stochastic
policy under an entropy-regularized objective, explicitly encouraging exploration.

It uses the same replay buffer configuration (\texttt{buffer\_size = 20{,}000}) and batch
size (\texttt{256}) as TD3 and DDPG for consistency. The entropy coefficient is set to
\texttt{ent\_coef = auto}, allowing the algorithm to automatically adjust the exploration–
exploitation trade-off during training.

As with other off-policy methods, \texttt{learning\_starts = 1000} ensures that early updates
are performed only after sufficient exploratory data has been collected.

Overall, on-policy and off-policy methods represent two complementary design choices in DRL.
On-policy algorithms rely on fresh, on-distribution data and typically offer more stable and
predictable learning dynamics, but at the cost of higher sample requirements. Off-policy
methods instead prioritize sample efficiency by reusing past experiences through replay buffers,
which is particularly beneficial in financial environments where data is expensive and sequential
correlations are strong. However, this comes with increased sensitivity to hyperparameters and
potential instability due to distributional shift in stored experiences. Evaluating both
families under carefully tuned configurations provides a more complete understanding of their
behavior in realistic trading scenarios.